\section{Lattice studies of the Wilson flow}

In QCD the perturbation expansion~\eqref{eq:E3}
is expected to be applicable
at small flow times only, where the smoothing range
$\sqrt{8t}$ is at most $0.3$ fm or so.
The properties of the Wilson flow at larger values of $t$ can however
be studied straightforwardly using the lattice formulation of the
theory and numerical simulations.

The simulations of the $\SUthree$ gauge theory
reported in this section mainly serve to clarify whether
$\langle E\rangle$ scales to the continuum limit as
suggested by perturbation theory. Along the way,
a new scale-setting method will be proposed based on the
observed properties of $\langle E\rangle$.

\begin{table}[ht]
\centering
\caption{Lattice parameters, statistics and reference flow time}
\label{tab:params}
\begin{tabular}{ccccc}
\hline
Lattice & $\beta$ & $a$\,[fm] & $N_{\mathrm{cnfg}}$ & $t_0/a^2$ \\
\hline
$48\times24^3$ & $5.96$ & $0.0999(4)$ & $100$ & $\phantom{0}2.7854(62)$ \\
$64\times32^3$ & $6.17$ & $0.0710(3)$ & $100$ & $\phantom{0}5.489(14)\phantom{0}$ \\
$96\times48^3$ & $6.42$ & $0.0498(3)$ & $100$ & $11.241(23)$ \\
\hline
\end{tabular}
\end{table}

\subsection{Simulation parameters}

Ensembles of representative gauge fields were
generated on three lattices, using the Wilson gauge action
and a combination of the well-known link-update algorithms.
The values of the lattice spacing quoted in table~\ref{tab:params}
derive from the results in lattice units for the Sommer
reference scale $r_0=0.5$ fm~\cite{SommerScale}
published by Guagnelli et al.~\cite{GuagnelliEtAl}.
At the chosen couplings $\beta=6/g_0^2$,
the spacings of the three lattices thus
decrease from roughly $0.1$ to $0.05$ fm by factors of $1/\sqrt{2}$.
Moreover, the lattice sizes in physical units are approximately
constant.

In order to safely suppress any residual statistical correlations of
the $N_{\mathrm{cnfg}}$ generated fields, the separation in simulation time
of the fields was taken to be at least $10$ times
the integrated autocorrelation time of the topological
charge. The definition of the latter on the lattice
is ambiguous to some extent, but
its autocorrelation time is largely independent
of the choices one makes and is known to increase very rapidly
when the lattice spacing is reduced~\cite{DelDebbioTauQ,StefanConference}.
In particular, already at $a=0.07$ fm
all other usual quantities of interest tend to be far less
correlated in simulation time.


\subsection{Observables}

For any given gauge field configuration $U(x,\mu)$, the flow equation~\eqref{eq:wflow}
can be integrated numerically up to the desired flow time and
one may then construct gauge-invariant local observables from
the gauge field at this time. In particular,
\begin{equation}
  E=2\sum_{p\in P_x}\Re\tr\{1-V_t(p)\}
  \label{eq:Elattice}
\end{equation}
is a possible definition of the density $E$ on the lattice,
$P_x$ being the set of unoriented plaquettes
with lower-left corner $x$.

\begin{figure}[ht]
\centering
\includegraphics[width=0.42\textwidth]{images/clover}
\caption{The anti-hermitian traceless
part of the average of the four plaquette Wilson loops
shown in this figure can be taken as the
definition of the field tensor $a^2G_{\mu\nu}(x)$
on the lattice. All loops are in the
$(\mu,\nu)$-plane, have the same orientation
and start and end at the point $x$.}
\label{fig:clover}
\end{figure}

Another more symmetric definition of $E$ is
obtained by introducing a lattice version of the
field tensor $G_{\mu\nu}(x)$ (see fig.~\ref{fig:clover}).
$E$ is then simply given by the continuum formula~\eqref{eq:Edens}.
Both definitions
are equally acceptable at this point,
since they respect all formal requirements
(locality and gauge invariance in particular) and since
they converge to
the correct expression in the classical continuum limit.

For the numerical integration of the flow equation~\eqref{eq:wflow},
any of the widely known
integration schemes can in principle be used
(see ref.~\cite{HairerEtAl}, for example).
The Euler integrators previously discussed in ref.~\cite{TrivMaps}
are particularly simple but also the least efficient ones.
Evidently, the integration errors should be much smaller than
the statistical errors of the calculated quantities.
A fairly simple and numerically stable
integrator that allows this condition to be easily met
is described in Appendix~\ref{app:C}.

\begin{figure}[ht]
\centering
\includegraphics[width=0.62\textwidth]{images/tsqE}
\caption{Simulation data obtained at $a=0.05$ fm for
$t^2\langle E\rangle$ as a function of the flow time $t$
(black line). Statistical errors are smaller than $0.3\%$
and therefore invisible on the scale of the figure.
The curve predicted by the perturbation series~\eqref{eq:E3}
and the known value~\eqref{eq:Lambdaparm} of the $\Lambda$-parameter
is also shown (grey band).}
\label{fig:tsqE}
\end{figure}

\subsection{Time dependence of $\langle E\rangle$}

To leading-order perturbation theory, the dimensionless combination
$t^2\langle E\rangle$ is a constant proportional to the gauge coupling.
At the next order, the scale invariance of the theory
is broken and $t^2\langle E\rangle$ develops a
non-trivial dependence on
the flow time $t$. Asymptotic freedom actually implies that
the combination slowly goes to zero in the limit $t\to0$.

Recalling the result
\begin{equation}
  \left.\Lambda\right|_{\Nf=0}=0.602(48)/r_0
  \label{eq:Lambdaparm}
\end{equation}
for the $\Lambda$-parameter in the $\MSbar$ scheme
obtained by the ALPHA collaboration~\cite{LambdaParm},
and using the four-loop evolution equation
for the running coupling $\alpha(q)$~\cite{RitbergenEtAl},
the perturbation series~\eqref{eq:E3} can be
evaluated at any value of the flow time
given in units of $r_0$.
The curve obtained in this way
is shown in fig.~\ref{fig:tsqE} together with the error band that derives
from the error of $\Lambda$ quoted in eq.~\eqref{eq:Lambdaparm}.

Perhaps somewhat fortuitously, the simulation results obtained
at $a=0.05$ fm accurately match the perturbative curve
over a significant range of $t$.
The symmetric definition of $E$ has here been used and
for clarity the data from only one lattice are shown
(as discussed below, the lattice-spacing effects are small
and the data from the other lattices would therefore lie nearly on
top of the line plotted in fig.~\ref{fig:tsqE}).

Beyond the perturbative regime, $t^2\langle E\rangle$ grows
roughly linearly with $t$, at least so within the
range covered by the simulation data.
The slowdown of the density
$\langle E\rangle$ from the
perturbative $1/t^2$ to a smoother $1/t$ behaviour
may perhaps be explained by noting that
the Wilson flow tends to drive the gauge field towards the
stationary points of the gauge action. In the vicinity
of these points of field space, the right-hand
side of the flow equation~\eqref{eq:wflow} is small and $E$
consequently changes only little with time.


\subsection{Lattice-spacing effects}

Perturbation theory suggests that
the density $\langle E\rangle$ scales to the continuum
limit like a physical quantity of dimension $4$.
The scaling behaviour of $\langle E\rangle$ can be
checked by introducing a reference scale $t_0$ through
the implicit equation
\begin{equation}
  \left\{t^2\langle E\rangle\right\}_{t=t_0}=0.3
  \label{eq:t0def}
\end{equation}
(see fig.~\ref{fig:tsqE}). If $\langle E\rangle$ is physical,
the dimensionless ratio $t_0/r_0^2$ must be independent
of the lattice spacing, up to corrections vanishing
proportionally to a power of $a$.

\begin{figure}[ht]
\centering
\includegraphics[width=0.62\textwidth]{images/t0r0}
\caption{Extrapolation of the dimensionless ratio $\sqrt{8t_0}/r_0$ of
reference scales to
the continuum limit (open data points). The black
data points were obtained using the symmetric definition of
$E$ and the grey ones using the expression~\eqref{eq:Elattice}.}
\label{fig:t0r0}
\end{figure}

The data plotted in fig.~\ref{fig:t0r0} clearly show that the ratio of
reference scales smoothly extrapolates to the continuum limit.
As may have been suspected, the lattice effects appear to
be of order $a^2$. They are at most a few percent on the
lattices considered and particularly small
if the symmetric definition of $E$ is employed.
At other points in time, the scaling violations behave similarly,
but tend to increase towards small $t$, where
the smoothing range $\sqrt{8t}$ is only 2 or 3 times larger than
the lattice spacing.


\subsection{Synthesis}

(a) \textit{Continuum limit}.
The numerical studies
of the Wilson flow strongly support the conjecture that
the gauge fields generated at positive flow time
are smooth renormalized fields (except for their gauge
degrees of freedom). In the case of QCD with a non-zero number
of sea quarks, similar studies however still need to be performed.
Additional confirmation in perturbation theory and through
the consideration of further observables is evidently
desirable.

\vspace{0.5ex}
\noindent
(b) \textit{Scale setting}.
The time $t_0$ defined through eq.~\eqref{eq:t0def} may serve as
a reference scale similar to the Sommer radius $r_0$.
With respect to the latter, $t_0$ has the advantage that
its computation does not require any fits or extrapolations.
Moreover, an ensemble of only 100 independent
representative field configurations allows
$t_0$ to be obtained
with a statistical precision of a small fraction of a percent
(for illustration, the values of $t_0/a^2$ computed
using the symmetric definition of $E$ are listed in table~\ref{tab:params}).

\vspace{0.5ex}
\noindent
(c) \textit{Universality}.
On the lattice the definition of
local gauge-invariant quantities like $E$ is not unique,
but the results obtained in this section indicate that
the differences become irrelevant in the continuum limit
(see fig.~\ref{fig:t0r0}).
This kind of universality is a consequence of the fact,
further elucidated in sect.~4, that
the gauge fields generated by the Wilson flow are smooth
on the scale of the lattice spacing.
The universality classes of the local fields composed from the
gauge field at fixed $t/t_0$
are therefore determined by the asymptotic behaviour of the fields
in the classical continuum limit.

